\usepackage[T1]{fontenc}
\usepackage[utf8]{inputenc}
\newcommand\hmmax{0} % Restricts heavy math alphabets
\newcommand\bmmax{0}
\usepackage{amsmath}
\usepackage{gensymb}
\usepackage{tabularx}
\usepackage{xltabular}
\usepackage{cancel}
\usepackage{geometry}
\setlength{\arrayrulewidth}{0.25mm}
\usepackage{array}
\setlength{\tabcolsep}{5pt}
\renewcommand{\arraystretch}{1.25}
\usepackage[output-decimal-marker={,},exponent-product=\cdot]{siunitx}
\usepackage{siunitx,upgreek,abbrevs}
\usepackage{enumitem}
\usepackage[french]{babel}
\usepackage[french]{datetime2}
\DeclareMathAlphabet{\mathbb}{U}{msb}{m}{n}
\usepackage{mathtools}
\usepackage{sagetex}
\usepackage[labelfont=bf,textfont=md,margin=1cm]{caption}
%\usepackage{bm}
\usepackage{nicematrix}
\usepackage{makecell}
\usepackage{pdfpages}
\usepackage{fancyhdr}
\usepackage{lastpage}
\usepackage{circuitikz}
\usepackage{pdfcomment}
\usepackage{comment}
\usepackage{pgfplots}
\usepackage{pgfplotstable}
\usepackage{subcaption}
\usepackage{changepage}
\usepackage{hyperref}
\usepackage{nicefrac}
\usepackage{newunicodechar}
\geometry{margin=0.65in}

\newcommand{\units}[1]{\,\unit{#1}}
\newcommand{\dii}{\mathop{\mathrm{d}\!}}
\newcommand{\dif}[3][]{\frac{\mathrm{d}^{#1}}{\mathrm{d}#3^{#1}}\left(#2\right)}
\newcommand{\difs}[3][]{\frac{\mathrm{d}^{#1}#2}{\mathrm{d}#3^{#1}}}
\newcommand{\pif}[3][]{\frac{\partial^{#1}}{\partial#3^{#1}}\left(#2\right)}
\newcommand{\pifs}[3][]{\frac{\partial^{#1}#2}{\partial#3^{#1}}}
\newcommand{\ddif}[3][]{\dfrac{\mathrm{d}^{#1}}{\mathrm{d}#3^{#1}}\left(#2\right)}
\newcommand{\ddifs}[3][]{\dfrac{\mathrm{d}^{#1}#2}{\mathrm{d}#3^{#1}}}
\newcommand{\dpif}[3][]{\dfrac{\partial^{#1}}{\partial#3^{#1}}\left(#2\right)}
\newcommand{\dpifs}[3][]{\dfrac{\partial^{#1}#2}{\partial#3^{#1}}}
\newcommand{\solution}{\textbf{\underline{\textsc{Solution}}}\\}
\newcommand{\abs}[1]{\left|#1\right|}
\newcommand{\prb}[1]{\left(#1\right)}
\newcommand{\brb}[1]{\left[#1\right]}
\newcommand{\vnorm}[1]{\left\lVert#1\right\rVert}
\newcommand{\sln}{\boxed{\textbf{Solution}}}
\let\oldforall\forall
\renewcommand{\forall}{\:\oldforall\:}
\newcommand{\firstpageheadrule}{\renewcommand{\headrulewidth}{0.4pt}}
\newcommand{\runningheadrule}{\renewcommand{\headrulewidth}{0.4pt}}
\newcommand{\firstpagefootrule}{\renewcommand{\footrulewidth}{0.4pt}}
\newcommand{\runningfootrule}{\renewcommand{\footrulewidth}{0.4pt}}
\renewcommand{\cellalign}{cc}
\newcommand{\fvspacing}{20pt}
\newunicodechar{°}{\degree}
\usepackage{xstring}

% Obligatory Gemini credit for this block
\ExplSyntaxOn
\newcommand{\autofig}[3][scale=1]{
	\StrSubstitute{#2}{.png}{}[\tempname]
	\edef\cleanname{\detokenize\expandafter{\tempname}}
	\bool_set_false:N \l_tmpa_bool
	\file_if_exist:nT { .\jobname.tex_eet } { \bool_set_true:N \l_tmpa_bool }
	\file_if_exist:nT { .\jobname .tex_et } { \bool_set_true:N \l_tmpa_bool }
	\bool_if:NTF \l_tmpa_bool
	{
		\pdftooltip{\includegraphics[#1]{#2}}{\cleanname}
	}
	{
		% No flag found: Regular image
		\includegraphics[#1]{#2}
	}
	\caption{#3.}
	\label{fig:\cleanname}
}
\ExplSyntaxOff

\newif\ifneedscolor
\needscolorfalse

\foreach \env in {gather, gather*, align, align*, equation, equation*} {
		\AddToHook{env/\env/begin}{\needscolortrue}
		\AddToHook{env/\env/end}{\needscolorfalse}
	}

\NewDocumentCommand{\qtys}{m m m}{%
	\ifneedscolor\sagestr{qty(#1,#2,"#3")}\else\sagestr{qty(#1,#2,"#3",False)}\fi
}
\NewDocumentCommand{\nums}{m m O{None} O{None}}{%
	\ifneedscolor\sagestr{num(#1,#2,True,#3,#4)}\else\sagestr{num(#1,#2,False,#3,#4)}\fi
}

\newcommand{\datai}[3]{\nums{#1}{#3}[False][#2]\pm\nums{#2}{#3}[True][None]}
\newcommand{\datail}[2]{\nums{#1[0]}{#2}[False][#1[1]]\pm\nums{#1[1]}{#2}[True][None]}
\newcommand{\datailt}[3]{%
	\ifnum#2<0
		\left[\datail{#1}{#2}\right]\!\units{#3}%
	\else
		\left(\datail{#1}{#2}\right)\!\units{#3}%
	\fi
}
\newcommand{\maxl}[1]{#1_\text{max.}}
\newcommand{\effl}[1]{#1_\text{eff.}}

\DeclareMathVersion{myques}
\SetSymbolFont{operators}{myques}{OT1}{cmss}{m}{n}
\SetSymbolFont{letters}  {myques}{OT1}{cmss}{m}{it}
\SetSymbolFont{symbols}  {myques}{OMS}{cmsy}{m}{n}
\SetMathAlphabet{\mathsf}{myques}{OT1}{cmss}{m}{n}
\SetMathAlphabet{\mathbf}{myques}{OT1}{cmss}{bx}{n}
\SetMathAlphabet{\mathrm}{myques}{OT1}{cmss}{m}{n}
\DeclareSymbolFont{sfgreek}{OML}{ntxsfmi}{m}{it}
\DeclareMathSymbol{\sfalpha}{\mathord}{sfgreek}{"0B}
\DeclareMathSymbol{\sfbeta}{\mathord}{sfgreek}{"0C}
\DeclareMathSymbol{\sfgamma}{\mathord}{sfgreek}{"0D}
\DeclareMathSymbol{\sfdelta}{\mathord}{sfgreek}{"0E}
\DeclareMathSymbol{\sfepsilon}{\mathord}{sfgreek}{"0F}
\DeclareMathSymbol{\sftheta}{\mathord}{sfgreek}{"12}
\DeclareMathSymbol{\sflambda}{\mathord}{sfgreek}{"15}
\DeclareMathSymbol{\sfmu}{\mathord}{sfgreek}{"16}
\DeclareMathSymbol{\sfpi}{\mathord}{sfgreek}{"19}
\DeclareMathSymbol{\sfrho}{\mathord}{sfgreek}{"1A}
\DeclareMathSymbol{\sfsigma}{\mathord}{sfgreek}{"1B}
\DeclareMathSymbol{\sfphi}{\mathord}{sfgreek}{"1E}
\DeclareMathSymbol{\sfomega}{\mathord}{sfgreek}{"21}
\DeclareMathSymbol{\sfGamma}{\mathord}{sfgreek}{"00}
\DeclareMathSymbol{\sfDelta}{\mathord}{sfgreek}{"01}
\DeclareMathSymbol{\sfTheta}{\mathord}{sfgreek}{"02}
\DeclareMathSymbol{\sfLambda}{\mathord}{sfgreek}{"03}
\DeclareMathSymbol{\sfXi}{\mathord}{sfgreek}{"04}
\DeclareMathSymbol{\sfPi}{\mathord}{sfgreek}{"05}
\DeclareMathSymbol{\sfSigma}{\mathord}{sfgreek}{"06}
\DeclareMathSymbol{\sfUpsilon}{\mathord}{sfgreek}{"07}
\DeclareMathSymbol{\sfPhi}{\mathord}{sfgreek}{"08}
\DeclareMathSymbol{\sfPsi}{\mathord}{sfgreek}{"09}
\DeclareMathSymbol{\sfOmega}{\mathord}{sfgreek}{"0A}
\DeclareMathSymbol{\sfnu}{\mathord}{sfgreek}{"17}
% 1. Declare CM integral symbol in the preamble
\DeclareSymbolFont{CMlargesymbols}{OMX}{cmex}{m}{n}
\DeclareMathSymbol{\CMintop}{\mathop}{CMlargesymbols}{"52}

\newcommand{\questionenum}[1]{%
	{\sffamily
		\mathversion{myques}%
		\mathcode`.="002E\relax
		\mathcode`,="602C\relax
		\let\alpha\sfalpha \let\beta\sfbeta \let\gamma\sfgamma
		\let\delta\sfdelta \let\epsilon\sfepsilon \let\theta\sftheta
		\let\lambda\sflambda \let\mu\sfmu \let\pi\sfpi
		\let\rho\sfrho \let\sigma\sfsigma \let\phi\sfphi
		\let\omega\sfomega \let\Gamma\sfGamma \let\Delta\sfDelta \let\Theta\sfTheta \let\nu\sfnu
		\let\Lambda\sfLambda \let\Xi\sfXi \let\Pi\sfPi
		\let\Sigma\sfSigma \let\Upsilon\sfUpsilon \let\Phi\sfPhi
		\let\Psi\sfPsi \let\Omega\sfOmega
		% 2. Save normal \intop and switch to \CMintop for the question
		\let\origintop\intop
		\let\intop\CMintop
		\sisetup{
			mode = match,
			propagate-math-font = true,
			reset-math-version = false,
			output-decimal-marker = {\textsf{,}}
		}%
		#1\setlength{\parskip}{0pt}%
		\par\nobreak
		\noindent\rule{\linewidth}{0.4pt}%
		\par\nobreak
		% 3. Restore normal \intop for \sln
		\noindent{\normalfont\mathversion{normal}\let\intop\origintop\sln}\quad%
	}%
}
\newcommand{\question}[1]{\questionenum{#1}}
\makeatletter
\@ifpackageloaded{MinionPro}{%
	% 1. Font encodings & symbol fonts from newtx
	\DeclareFontEncoding{LMS}{}{}%
	\DeclareFontSubstitution{LMS}{ntxsy}{m}{n}%

	\DeclareSymbolFont{NTXsymbols}{LMS}{ntxsy}{m}{n}%
	\SetSymbolFont{NTXsymbols}{bold}{LMS}{ntxsy}{b}{n}%

	\DeclareSymbolFont{NTXletters}{OML}{ntxmi}{m}{it}%
	\SetSymbolFont{NTXletters}{bold}{OML}{ntxmi}{b}{it}%

	\DeclareSymbolFont{NTXoperators}{OT1}{ntxtlf}{m}{n}%
	\SetSymbolFont{NTXoperators}{bold}{OT1}{ntxtlf}{b}{n}%

	% Scaling helper for math relations (default 0.9)
	\newcommand{\NTXscale}[2][0.9]{%
		\mathrel{%
			\mathchoice
			{\scalebox{#1}{$\displaystyle#2$}}%
			{\scalebox{#1}{$\textstyle#2$}}%
			{\scalebox{#1}{$\scriptstyle#2$}}%
			{\scalebox{#1}{$\scriptscriptstyle#2$}}%
		}%
	}%

	% 2. Basic comparisons: =, <, >
	\DeclareMathSymbol{\NTXeq}{\mathrel}{NTXoperators}{"3D}%
	\DeclareMathSymbol{=}{\mathrel}{NTXoperators}{"3D}%

	\DeclareMathSymbol{\NTXlt}{\mathrel}{NTXletters}{"3C}%
	\DeclareMathSymbol{<}{\mathrel}{NTXletters}{"3C}%

	\DeclareMathSymbol{\NTXgt}{\mathrel}{NTXletters}{"3E}%
	\DeclareMathSymbol{>}{\mathrel}{NTXletters}{"3E}%

	% 3. Inequality symbols: \neq/\ne, \le/\leq, \ge/\geq
	\DeclareMathSymbol{\NTXleq}{\mathrel}{NTXsymbols}{"14}%
	\DeclareMathSymbol{\NTXgeq}{\mathrel}{NTXsymbols}{"15}%

	\newcommand{\@scaledneq@glyph}[2]{%
		\scalebox{#1}{$#2\mathrel{\ooalign{\hfil$#2\NTXeq$\hfil\cr\hfil$#2/$\hfil}}$}%
	}%
	\newcommand{\scaledneq}[1][0.9]{%
		\mathrel{%
			\mathchoice
			{\@scaledneq@glyph{#1}{\displaystyle}}%
			{\@scaledneq@glyph{#1}{\textstyle}}%
			{\@scaledneq@glyph{#1}{\scriptstyle}}%
			{\@scaledneq@glyph{#1}{\scriptscriptstyle}}%
		}%
	}%

	% 4. Equivalence & set relations from NTXsymbols (LMS)
	\DeclareMathSymbol{\NTXapprox}{\mathrel}{NTXsymbols}{"19}%
	\DeclareMathSymbol{\NTXsimeq}{\mathrel}{NTXsymbols}{"27}%
	\DeclareMathSymbol{\NTXsim}{\mathrel}{NTXsymbols}{"18}%
	\DeclareMathSymbol{\NTXequiv}{\mathrel}{NTXsymbols}{"11}%
	\DeclareMathSymbol{\NTXsubset}{\mathrel}{NTXsymbols}{"1A}%
	\DeclareMathSymbol{\NTXsupset}{\mathrel}{NTXsymbols}{"1B}%
	\DeclareMathSymbol{\NTXsubseteq}{\mathrel}{NTXsymbols}{"12}%
	\DeclareMathSymbol{\NTXsupseteq}{\mathrel}{NTXsymbols}{"13}%

	% 5. Apply replacements at \AtBeginDocument so MinionPro/MnSymbol/amsmath don't clobber them
	\AtBeginDocument{%
		\renewcommand{\approx}{\NTXscale{\NTXapprox}}%
		\renewcommand{\simeq}{\NTXscale{\NTXsimeq}}%
		\renewcommand{\sim}{\NTXscale{\NTXsim}}%
		\renewcommand{\equiv}{\NTXscale{\NTXequiv}}%
		\renewcommand{\subset}{\NTXscale{\NTXsubset}}%
		\renewcommand{\supset}{\NTXscale{\NTXsupset}}%
		\renewcommand{\subseteq}{\NTXscale{\NTXsubseteq}}%
		\renewcommand{\supseteq}{\NTXscale{\NTXsupseteq}}%
		\renewcommand{\le}{\NTXscale{\NTXleq}}%
		\renewcommand{\leq}{\NTXscale{\NTXleq}}%
		\renewcommand{\ge}{\NTXscale{\NTXgeq}}%
		\renewcommand{\geq}{\NTXscale{\NTXgeq}}%
		\renewcommand{\neq}{\scaledneq}%
		\renewcommand{\ne}{\scaledneq}%
	}%
}{}
\makeatother

\newcommand{\sfig}[2][1]{\begin{figure}[h!]\centering\linespread{1}\scalebox{#1}{#2}\end{figure}}
\newcommand{\parallelsum}{\mathbin{\!/\mkern-5mu/\!}}
\newcommand{\invs}{s^{-1}}
\newcommand{\cond}[2]{\text{cond}_{#1}\!\left(#2\right)}
\newcommand{\nfrac}{\nicefrac}

\usetikzlibrary{babel}
\usetikzlibrary{arrows.meta}
\tikzset{>={Latex[width=2.5mm,length=2.5mm]}}
\linespread{1.2}
\NiceMatrixOptions{cell-space-top-limit=4pt,cell-space-bottom-limit=2pt}

\makeatletter
\renewcommand*\env@matrix[1][*\c@MaxMatrixCols c]{%
	\hskip -\arraycolsep
	\let\@ifnextchar\new@ifnextchar
	\array{#1}}
\makeatother

\setcounter{secnumdepth}{5}
\setcounter{tocdepth}{5}

\makeatletter
\newcommand\subsubsubsection{\@startsection{paragraph}{4}{\z@}{-2.5ex\@plus -1ex \@minus -.25ex}{1.25ex \@plus .25ex}{\normalfont\normalsize\itshape}}
\makeatother

\pgfplotsset{
	filter discard warning=false
}
